\documentclass[10pt,a4paper]{amsart}
\usepackage[margin=19mm]{geometry}
\usepackage{amsmath,amssymb,mathtools}
\usepackage[scr=rsfs,cal=euler]{mathalfa}
\usepackage{xfrac}
\usepackage[unicode,hidelinks,bookmarks=false]{hyperref}
\newcommand{\ie}{i.e.\ }
\newcommand{\cf}{cf.\ }
\newcommand{\resp}{respectively\ }
\newcommand{\Subsection}{\subsection}
\providecommand{\nobreakdash}{\nobreak\hbox{-}}
\setlength{\emergencystretch}{2em}
\multlinegap0pt
\usepackage{amssymb}
\usepackage{bm}
\usepackage{mathtools}
\usepackage[shortlabels]{enumitem}
\usepackage{float}
\usepackage{xcolor}
\usepackage{tikz}
\usepackage{tcolorbox}
\newtheorem{lemma}{Lemma}
\newcommand{\qbin}[2]{\genfrac{[}{]}{0pt}{}{#1}{#2}_{q}}
\title{Formalized $q$-series: The Rogers--Ramanujan Identities and Beyond}
\author{Kenny Lau, Seewoo Lee, Ken Ono}
\date{Source: arXiv:2607.01544v3 (CC BY 4.0)}
\begin{document}
\maketitle

\noindent\emph{Source and licence.} Formalized $q$-series: The Rogers--Ramanujan Identities and Beyond. Version arXiv:2607.01544v3 (CC BY 4.0), distributed by arXiv under the Creative Commons Attribution 4.0 International licence, http://creativecommons.org/licenses/by/4.0/, as recorded in the arXiv licence metadata for this exact version and shown on the abstract page https://arxiv.org/abs/2607.01544v3. Full licence notice: \texttt{licenses/arxiv-2607.01544-CC-BY-4.0.txt}.\par\smallskip

\noindent\textit{Editorial scope.} Counted unit: the article's lemma lem:binomial\_cauchy (source0.tex lines 831--838). The article's complete original proof of it is delivered with it (source0.tex lines 899--924, one \texttt{proof} environment of 26 lines, which the article places away from the statement). No mathematics is written, repaired or re-derived: the statement and the proof are the article's own lines, in the article's own notation, and no retained line is substituted or re-flowed.  Retained uncounted as the source-local context the proof needs: lines 859--865.  Nothing was omitted from any retained span, so the package carries no omission record.  No retained line contains a citation, so the wrapper carries no bibliography and no bibliography entry.  No retained line contains a figure environment, an image-inclusion command or drawing code, so no image is delivered and the package declares no asset dependency.  Editorial formatting only, in addition: an AMS \texttt{amsart} preamble that restates the article's own declarations and macros used by the retained lines, namely the environments \texttt{lemma} on separate counters, together with the standard packages those lines need.  The retained environments print in this document as Lemma 1 in order of appearance, whereas the article prints the same items with the numbering of its own full document.  The complete native source of the article as distributed in the arXiv e-print for this exact version is bundled unchanged as \texttt{sources/204263/source0.tex}.
\begin{lemma}
    \label{lem:binomial_mul}
    For $n, k, s$ with $s \le k$, we have
    \begin{equation}
        \qbin{n}{k} \qbin{k}{s} = \qbin{n}{s} \qbin{n-s}{k-s}
    \end{equation}
\end{lemma}

\begin{lemma}[$q$-Cauchy identity]
    \label{lem:binomial_cauchy}
    \begin{equation}
        (uv)_n = \sum_{k=0}^{n} \begin{bmatrix}
            n \\ k
        \end{bmatrix}_q (u)_k v^k (v)_{n-k}
    \end{equation}
\end{lemma}

\begin{proof}
    Applying Lemma \ref{lem:binomial_cauchy} to $u = bq^k$ and $v = ac$, we can write
    \begin{equation}
        \label{eqn:qPS-lemma}
        (abcq^k)_{n-k} = \sum_{r=0}^{n-k} \begin{bmatrix}
            n-k \\ r
        \end{bmatrix}_q (bq^k)_r (ac)^r (ac)_{n-k-r}
    \end{equation}
    and the right hand side becomes
    \begin{align*}
        &\sum_{k=0}^{n} \sum_{r=0}^{n-k} \qbin{n}{k} \qbin{n-k}{r} (a)_k (b)_k c^k (c)_{n-k} (bq^k)_r (ac)^r (ac)_{n-k-r}\\
        &=\sum_{k=0}^{n} \sum_{r=0}^{n-k} \qbin{n}{k+r}\qbin{k+r}{k} (a)_k (b)_{k+r} a^r c^{k+r} (c)_{n-k} (ac)_{n-k-r} \\
        &= \sum_{m=0}^{n} \qbin{n}{m} (b)_m c^m (c)_{n-m} (ac)_{n-m} \sum_{k=0}^{m} \qbin{m}{k} (a)_k a^{m-k} (cq^{n-m})_{m-k}.
    \end{align*}
    Here we rearranged the double sums, where $m = k + r$ for the second equality.
    We also used Lemma \eqref{lem:binomial_mul}, $(b)_k (bq^k)_r = (b)_{k+r}$ and $(c)_{n-k} = (c)_{n-m} (cq^{n-m})_{m-k}$.
    By Lemma \ref{lem:binomial_cauchy} to $u = cq^{n-m}$ and $v = a$, the inner sum becomes
    $$
        \sum_{k=0}^{m} \qbin{m}{k} (a)_k a^{m-k} (cq^{n-m})_{m-k} = (acq^{n-m})_m.
    $$
    Plugging this into the last sum and using $(ac)_{n-m}(acq^{n-m})_m = (ac)_n$ gives
    \begin{align*}
        (ac)_n \sum_{m=0}^{n} \qbin{n}{m} (b)_m c^m (c)_{n-m} = (ac)_n (bc)_n
    \end{align*}
    by Lemma \ref{lem:binomial_cauchy} again.
\end{proof}

\end{document}
